\chapter{Lezione 21 --- Bande di Gershgorin, criterio di Nyquist per sistemi MIMO, SVD e minimi quadrati}

\section{Sistemi MIMO e dominanza diagonale}

Consideriamo un sistema dinamico MIMO descritto in frequenza da
\[
X(s)=G(s)U(s),
\qquad
X(s)\in\mathbb{C}^m,\quad U(s)\in\mathbb{C}^m,
\]
dove
\[
G(s)\in\mathbb{C}^{m\times m}.
\]

Se il sistema fosse completamente disaccoppiato, cioè con matrice di trasferimento diagonale,
\[
G(s)=
\begin{bmatrix}
g_{11}(s) & 0 & \cdots & 0\\
0 & g_{22}(s) & \cdots & 0\\
\vdots & \vdots & \ddots & \vdots\\
0 & 0 & \cdots & g_{mm}(s)
\end{bmatrix},
\]
potremmo studiare ogni canale separatamente, analizzando ciascun elemento diagonale.

Se invece il sistema non è diagonale ma \emph{a dominanza diagonale}, allora gli elementi fuori diagonale possono essere interpretati come perturbazioni rispetto ai termini diagonali dominanti.

\medskip

\noindent
\textbf{Dominanza diagonale per righe.}
Una matrice \(A\in\mathbb{R}^{m\times m}\) si dice diagonalmente dominante per righe se, per ogni \(i=1,\dots,m\),
\[
|a_{ii}| \gg \sum_{\substack{j=1\\j\neq i}}^{m}|a_{ij}|.
\]

\noindent
\textbf{Dominanza diagonale per colonne.}
Analogamente, si può richiedere che, per ogni \(i=1,\dots,m\),
\[
|a_{ii}| \gg \sum_{\substack{j=1\\j\neq i}}^{m}|a_{ji}|.
\]

Nel seguito usiamo la convenzione degli appunti, cioè quella associata alle \emph{colonne}.

\section{Teorema di Gershgorin per matrici di trasferimento}

Per ogni frequenza \(\omega\), la matrice
\[
G(j\omega)
\]
è una matrice complessa \(m\times m\). Possiamo quindi applicare il teorema di Gershgorin.

\begin{theorem}[Gershgorin]
Sia \(G(j\omega)\in\mathbb{C}^{m\times m}\). Allora tutti gli autovalori di \(G(j\omega)\) appartengono all'unione dei cerchi del piano complesso centrati nei termini diagonali \(g_{ii}(j\omega)\) e di raggio
\[
r_i(j\omega)=\sum_{\substack{j=1\\j\neq i}}^{m}|g_{ji}(j\omega)|,
\qquad i=1,\dots,m.
\]
\end{theorem}

I cerchi di raggio \(r_i(j\omega)\) sono detti \textbf{cerchi di Gershgorin}.  
Al variare di \(\omega\), il centro \(g_{ii}(j\omega)\) si muove lungo il diagramma polare del termine diagonale, mentre il raggio cambia con la frequenza.

L'inviluppo di tali cerchi al variare di \(\omega\) genera le cosiddette \textbf{bande di Gershgorin}.

\section{Criterio di Nyquist formulato sulle bande di Gershgorin}

Per un sistema MIMO diagonalmente dominante, si può usare il criterio di Nyquist non direttamente sugli autovalori della matrice, ma sulle bande di Gershgorin.

L'idea è la seguente:

\begin{itemize}
    \item per ogni \(\omega\), gli autovalori di \(G(j\omega)\) stanno dentro i cerchi di Gershgorin;
    \item quindi, al variare della frequenza, le traiettorie degli autovalori sono contenute dentro le bande di Gershgorin;
    \item se tali bande \emph{non avvolgono} il punto critico \(-1\), allora anche il sistema accoppiato risulta stabile in ciclo chiuso.
\end{itemize}

Questa è una \textbf{condizione sufficiente} di stabilità del sistema in ciclo chiuso.

\begin{figure}[h!]
\centering
\begin{tikzpicture}[scale=1.05,>=Latex]
    % axes
    \draw[->] (-3.8,0)--(3.8,0) node[right] {$\Re$};
    \draw[->] (0,-2.8)--(0,2.8) node[above] {$\Im$};

    % point -1
    \fill (-2,0) circle (1.4pt);
    \node[below] at (-2,0) {$-1$};

    % Nyquist-like curve
    \draw[very thick,green!50!black]
        (2.0,0.1)
        .. controls (1.0,-0.2) and (0.6,-1.7) ..
        (-0.8,-1.9)
        .. controls (-1.6,-1.8) and (-1.8,-0.5) ..
        (-1.2,-0.2)
        .. controls (-0.7,0.0) and (0.4,0.15) ..
        (1.8,0.1);

    % representative centers
    \coordinate (c1) at (1.3,0.05);
    \coordinate (c2) at (0.4,-0.7);
    \coordinate (c3) at (-0.5,-1.5);
    \coordinate (c4) at (-1.1,-1.2);
    \coordinate (c5) at (-1.25,-0.45);
    \coordinate (c6) at (-0.4,-0.15);

    % Gershgorin circles
    \draw[red!75!black] (c1) circle (0.32);
    \draw[red!75!black] (c2) circle (0.42);
    \draw[red!75!black] (c3) circle (0.35);
    \draw[red!75!black] (c4) circle (0.28);
    \draw[red!75!black] (c5) circle (0.22);
    \draw[red!75!black] (c6) circle (0.18);

    \node[green!50!black,right] at (1.7,-0.5) {$g_{ii}(j\omega)$};
    \node[red!75!black,left] at (-2.9,-1.7) {cerchi di Gershgorin};
\end{tikzpicture}
\caption{Idea qualitativa: il centro di ciascun cerchio si muove sul diagramma polare di \(g_{ii}(j\omega)\), mentre il raggio è dato dalla somma dei moduli dei termini fuori diagonale associati.}
\end{figure}

\section{Esempio qualitativo}

Negli appunti compare un esempio in cui si prende un termine diagonale del tipo
\[
g_{11}(s)=\frac{10}{1+\frac{\sqrt{2}}{4}s+\frac{1}{16}s^2}
\]
e un termine fuori diagonale del tipo
\[
g_{21}(s)=\frac{5}{s+2}.
\]

In tal caso il centro dei cerchi di Gershgorin è il diagramma polare di \(g_{11}(j\omega)\), mentre il raggio vale
\[
r_1(j\omega)=|g_{21}(j\omega)|=\left|\frac{5}{2+j\omega}\right|
=\frac{5}{\sqrt{4+\omega^2}}.
\]

Quindi, ad esempio,
\[
r_1(0)=\frac{5}{2}=2.5,
\qquad
r_1(10)=\frac{5}{\sqrt{104}}\approx 0.49.
\]

Per \(\omega=0\), se \(g_{11}(0)=10\), il cerchio iniziale taglia l'asse reale nel punto
\[
10-\frac{5}{2}=7.5.
\]

Questo spiega il disegno negli appunti: a bassa frequenza i cerchi sono più grandi, poi diventano progressivamente più piccoli all'aumentare della frequenza.

\section{SVD: decomposizione ai valori singolari}

Sia ora
\[
A\in\mathbb{R}^{n\times m},
\qquad n\ge m.
\]
Allora esistono matrici ortogonali
\[
U\in\mathbb{R}^{n\times n},
\qquad
V\in\mathbb{R}^{m\times m},
\]
e una matrice diagonale ``rettangolare'' \(\Sigma\in\mathbb{R}^{n\times m}\) tali che
\[
A=U\Sigma V^T,
\]
con
\[
\Sigma=
\begin{bmatrix}
\widetilde\Sigma\\
0
\end{bmatrix},
\qquad
\widetilde\Sigma=
\operatorname{diag}(\sigma_1,\dots,\sigma_m),
\qquad
\sigma_1\ge \sigma_2\ge \cdots \ge \sigma_m\ge 0.
\]

I numeri \(\sigma_1,\dots,\sigma_m\) sono i \textbf{valori singolari} di \(A\).

Se
\[
\sigma_m>0,
\]
allora \(A\) ha rango pieno.

\subsection{Interpretazione dei valori singolari}

I valori singolari forniscono una misura di quanto una matrice sia vicina alla singolarità:

\begin{itemize}
    \item se tutti i valori singolari sono ``lontani da zero'', la matrice è ben condizionata;
    \item se almeno uno è molto piccolo, la matrice è vicina alla non invertibilità.
\end{itemize}

In particolare, il rapporto
\[
\frac{\sigma_1}{\sigma_m}
\]
misura il condizionamento della matrice.

\section{Problema dei minimi quadrati}

Consideriamo il problema
\[
Ax \approx y,
\qquad
A\in\mathbb{R}^{n\times m},
\quad
x\in\mathbb{R}^m,
\quad
y\in\mathbb{R}^n,
\]
con \(n\ge m\).

In generale il sistema può essere sovradeterminato e non avere soluzione esatta. Si cerca allora \(x\) che minimizza lo scostamento tra \(Ax\) e \(y\), cioè
\[
\min_x \frac12\|Ax-y\|^2.
\]

Nel caso di errore gaussiano, questa scelta coincide con la stima ottima nel senso dei minimi quadrati.

\begin{figure}[h!]
\centering
\begin{tikzpicture}[scale=0.95,>=Latex]
    \draw[->] (0,0)--(6.5,0) node[right] {$x$};
    \draw[->] (0,0)--(0,4.8) node[above] {$y$};

    % fitted line
    \draw[thick,teal!70!black] (0.5,0.9)--(6.0,4.2);
    \node[teal!70!black,right] at (5.2,4.0) {$y^\star=Kx$};

    % points
    \foreach \p/\q in {
        0.8/1.4,
        1.3/1.1,
        2.0/1.8,
        2.7/2.4,
        3.1/2.2,
        3.8/3.1,
        4.3/2.9,
        5.0/3.8,
        5.6/3.6
    }{
        \fill (\p,\q) circle (1.5pt);
    }

    % residuals
    \draw[densely dashed] (0.8,1.4)--($(0.8,0.9)!0.12!(6.0,4.2)$);
    \draw[densely dashed] (2.7,2.4)--($(2.7,0.9)!0.42!(6.0,4.2)$);
    \draw[densely dashed] (4.3,2.9)--($(4.3,0.9)!0.70!(6.0,4.2)$);

    \node[below right] at (3.0,1.45) {residui};
\end{tikzpicture}
\caption{Interpretazione geometrica del problema dei minimi quadrati: si cerca la retta/modello che minimizza gli scarti.}
\end{figure}

\subsection{Equazioni normali}

Scriviamo il funzionale:
\[
J(x)=\frac12(Ax-y)^T(Ax-y)
=\frac12\left(x^TA^TAx+y^Ty-2x^TA^Ty\right).
\]

La condizione di ottimo del primo ordine è
\[
\nabla_x J(x)=0
\quad\Longrightarrow\quad
A^TAx-A^Ty=0.
\]

Quindi
\[
\boxed{
A^TAx=A^Ty.
}
\]

Se \(A^TA\) è invertibile, la soluzione dei minimi quadrati è
\[
\boxed{
x=(A^TA)^{-1}A^Ty.
}
\]

La matrice
\[
\boxed{
A^+=(A^TA)^{-1}A^T
}
\]
è la \textbf{pseudo-inversa} di \(A\) (nel caso di rango pieno sulle colonne).

\section{Pseudo-inversa tramite SVD}

Se
\[
A=U\Sigma V^T,
\]
allora
\[
A^TA = V\Sigma^T\Sigma V^T
      = V\widetilde\Sigma^2V^T.
\]

Perciò gli autovalori di \(A^TA\) sono
\[
\sigma_1^2,\dots,\sigma_m^2,
\]
cioè \textbf{i quadrati dei valori singolari} di \(A\).

La pseudo-inversa si può allora scrivere come
\[
\boxed{
A^+=V\widetilde\Sigma^{-1}U^T.
}
\]

Quindi la soluzione dei minimi quadrati diventa
\[
\boxed{
x=A^+y=V\widetilde\Sigma^{-1}U^Ty.
}
\]

\subsection{Interpretazione in coordinate ortonormali}

Se definiamo
\[
\theta = U^Ty,
\qquad
z = V^Tx,
\]
allora il problema si diagonalizza e diventa, componente per componente,
\[
z_i = \frac{\theta_i}{\sigma_i},
\qquad i=1,\dots,m.
\]

Questa lettura è molto utile:
\begin{itemize}
    \item \(U^Ty\) sono le coordinate di \(y\) nella base delle colonne di \(U\);
    \item \(V^Tx\) sono le coordinate di \(x\) nella base delle colonne di \(V\);
    \item il passaggio da \(\theta_i\) a \(z_i\) consiste nel dividere per il valore singolare \(\sigma_i\).
\end{itemize}

Se un \(\sigma_i\) è molto piccolo, il corrispondente coefficiente viene fortemente amplificato: questo spiega la sensibilità numerica del problema.

\section{Condizionamento numerico}

Supponiamo ora di avere una matrice quadrata invertibile
\[
Fx=y.
\]

Se il dato \(y\) è perturbato,
\[
y=y^\star+\delta y,
\]
allora anche la soluzione sarà perturbata:
\[
x=x^\star+\delta x,
\qquad
\delta x = F^{-1}\delta y.
\]

L'errore sulla soluzione dipende quindi da \(F^{-1}\).  
Una misura della sensibilità del problema è il \textbf{numero di condizionamento}
\[
\kappa(F)=\|F\|\,\|F^{-1}\|.
\]

Nel caso della norma \(2\),
\[
\boxed{
\kappa_2(F)=\frac{\sigma_{\max}(F)}{\sigma_{\min}(F)}.
}
\]

Se \(\kappa_2(F)\approx 1\), il problema è ben condizionato; se invece \(\kappa_2(F)\gg 1\), piccoli errori nei dati possono produrre grandi errori nella soluzione.

Nel caso della matrice \(A^TA\),
\[
\kappa_2(A^TA)=\frac{\sigma_1^2}{\sigma_m^2}=\kappa_2(A)^2.
\]

Questa è la ragione per cui le equazioni normali peggiorano il condizionamento.

\section{Sintesi della lezione}

In questa lezione abbiamo visto due blocchi di argomenti.

\subsection*{1. Stabilità di sistemi MIMO mediante Gershgorin e Nyquist}
\begin{itemize}
    \item Se un sistema MIMO è diagonalmente dominante, si può approssimare il comportamento considerando la diagonale come parte principale e i termini fuori diagonale come perturbazioni.
    \item Per ogni frequenza \(\omega\), gli autovalori di \(G(j\omega)\) stanno nei cerchi di Gershgorin.
    \item L'inviluppo dei cerchi al variare di \(\omega\) costituisce le bande di Gershgorin.
    \item Se tali bande non circondano il punto critico \(-1\), allora si ha una condizione sufficiente di stabilità in ciclo chiuso.
\end{itemize}

\subsection*{2. SVD, minimi quadrati e pseudo-inversa}
\begin{itemize}
    \item La SVD scrive una matrice come
    \[
    A=U\Sigma V^T.
    \]
    \item I valori singolari misurano quanto la matrice sia vicina alla singolarità.
    \item Il problema dei minimi quadrati
    \[
    \min_x \frac12\|Ax-y\|^2
    \]
    porta alle equazioni normali
    \[
    A^TAx=A^Ty.
    \]
    \item La soluzione è
    \[
    x=A^+y,
    \qquad
    A^+=(A^TA)^{-1}A^T = V\widetilde\Sigma^{-1}U^T.
    \]
    \item Il condizionamento dipende dal rapporto tra valore singolare massimo e minimo.
\end{itemize}